% Post-CP26-paper-freeze candidate. No theta nodality or monodromy conclusion.
\subsection{Odd polarized norms on every two-rank-zero Jacobian}

\begin{lemma}[Every odd multiplier on the same supersingular surface]
\label{lem:root-f0-all-odd-polarized-norms-candidate}
\status{Candidate; direct proof supplied, independent audit pending.}
Let $(A,\Theta)$ be an indecomposable principally polarized abelian
surface over an algebraically closed field of characteristic two,
with two-rank zero. For every positive odd integer $n$ there is
an etale self-isogeny $\psi_n:A\to A$ such that
\[
 \psi_n^\dagger\psi_n=[n],\qquad
 \deg\psi_n=n^2,\qquad \psi_n^*\Theta\equiv n\Theta.
\]
No assertion about the singularities of $\psi_n(\Theta)$ is made.
\end{lemma}
\begin{proof}
The complete quaternion-parameter proof in
Lemma~\ref{lem:uniform-seven-interpolation-candidate} supplies
automorphisms $g,h$ of the same genus-two Jacobian satisfying
$g^2=h^2=-1$ and $gh=-hg$. Each preserves the principal
polarization, so Rosati sends $g,h,gh$ to their negatives.

By the four-square theorem, write
$n=a^2+b^2+c^2+d^2$ with integers $a,b,c,d$, and set
$\psi_n=a+bg+ch+dgh$ in the integral endomorphism ring.
Quaternion multiplication gives
\[
 (a-bg-ch-dgh)(a+bg+ch+dgh)
                =[a^2+b^2+c^2+d^2]=[n].
\]
Thus $\psi_n$ is an isogeny and its pullback polarization
is $n\lambda_\Theta$. Its degree has square
$\deg[n]=n^4$, hence equals $n^2$. Its kernel lies in
$A[n]$, which is etale because $n$ is odd; hence so is
$\psi_n$.

We specify the arithmetic input rather than infer it from a
finite search. Hardy and Wright, fourth edition, Theorems~369
and~370, \S20.5, pp.~302--303, give the complete descent
proof; their Theorem~87, p.~70, supplies its initial congruence.
Both statements and their complete relevant proofs were read.
Here is the argument in quaternion notation. Euler's norm identity
makes products of sums of four squares representable, so it suffices
to treat primes. The prime two is immediate. For an odd prime $p$,
the sets of residues $x^2$ and $-1-y^2$, with
$0\le x,y\le(p-1)/2$, each have $(p+1)/2$ elements and intersect.
Consequently $mp=x^2+y^2+1$ for some $0<m<p$.
Choose the least positive such multiplier for a representation
$mp=x_0^2+x_1^2+x_2^2+x_3^2$.
If $m$ is even, pair coordinates of the same parity and replace
each pair by its half-sum and half-difference. This represents
$(m/2)p$, contrary to minimality. Thus, unless $m=1$, it is
odd and at least three.
Choose integers $y_i\equiv x_i\pmod m$ with $|y_i|<m/2$.
They cannot all vanish, since $m^2\mid mp$ would imply
$m\mid p$, whereas $1<m<p$. Their squared sum is $m m'$
with $0<m'<m$. For the integral quaternions
$x=x_0+x_1i+x_2j+x_3ij$ and $y=y_0+y_1i+y_2j+y_3ij$,
every coefficient of $x\overline y$ is divisible by $m$:
modulo $m$ that product equals $x\overline x=mp$.
The integral quaternion $x\overline y/m$ has norm $pm'$,
contradicting minimality. Hence $m=1$, proving the required
prime case and then the four-square theorem.
\end{proof}

\begin{remark}[The remaining geometric condition]
\label{rem:root-f0-all-odd-polarized-norms-scope-candidate}
\status{Candidate scope statement; no new monodromy degree.}
This construction removes an isogeny-existence question only.
The earlier good-theta norm-transfer theorem still requires
a nodal individual theta image, which has not been proved for
these arbitrary multipliers. A proposed transfer through a
boundary of three generic $|3\Theta|$ curves remains a separate
geometric and representation-theoretic obligation. If that
transfer is established for $|9\Theta|$, the present lemma
would supply every odd degree divisible by three on the same
fixed two-rank-zero surface. This conditional consequence is
not counted as completed progress.
\end{remark}
